\section{La misma campaña a dos escalas}
\label{sec:escalas}

Esta sección es el objeto del informe. La campaña que se ejecuta hoy está acotada
por la potencia de una máquina de cómputo; la que el centro podría ejecutar sobre
cinco máquinas es la misma campaña con una rejilla más ancha. Lo que sigue separa
lo que cambia de lo que no.

\subsection{Lo que no cambia con la escala}

Nada de lo definido en las secciones 2 a 7. En concreto:

\begin{itemize}[leftmargin=1.4em,itemsep=0.25em]
  \item \textbf{Los diez nodos} y su orden. Son propiedades del método, no del
        presupuesto.
  \item \textbf{Las definiciones}, incluida la de nivel, que es la que impide que
        una comparación parezca de un campo cuando varían dos.
  \item \textbf{El marco y su sello.} Un sello no es más barato ni más caro.
  \item \textbf{La unidad de reporte.} Nueve paneles, nunca agregados. Más cómputo
        no produce más proteínas, de modo que el suelo de población por celda es
        exactamente el mismo en las dos escalas.
  \item \textbf{La regla de cierre} y el estadístico que decide.
\end{itemize}

El punto sobre la población merece énfasis porque es contraintuitivo:
\textbf{la escala de cómputo no compra resolución estadística}. El MDE de
\cod{LK:CCO} es 0.0234 porque esa celda tiene 821 proteínas emparejadas, y tendrá
821 con una máquina y con cincuenta. Solo una ventana distinta, o más queries,
cambian eso.

\subsection{Lo que sí cambia}

Dos cosas, y solo dos.

\subsubsection*{La anchura de la rejilla}

Con poco cómputo la rejilla se explora en \textbf{cascada}: se decide un nodo, se
fija su ganador o su haz, y se pasa al siguiente con los demás campos quietos. Con
mucho cómputo se puede explorar \textbf{cruzada}: el producto cartesiano de varios
nodos a la vez.

La diferencia importa porque la cascada supone que las decisiones no interactúan, y
esa suposición es falsa en general. Si el mejor sustrato bajo una política de
donantes permisiva no es el mejor bajo una restrictiva, la cascada lo pierde y el
cruce lo encuentra. \textbf{Lo que el centro compraría con cinco máquinas es
precisamente la capacidad de detectar interacciones entre nodos}, que es la única
cosa que la versión acotada no puede ver por construcción.

\begin{center}
\begin{tabularx}{\linewidth}{@{}L{3.4cm} X X@{}}
\toprule
 & cascada (hoy) & cruce (a escala)\\
\midrule
arms de predicción & 25 del sustrato, luego 4 de política: \textbf{29}. Los cortes no añaden arms, truncan una lista ya recuperada & $22 \times 16 = \mathbf{352}$\\
supuesto           & los nodos no interactúan & ninguno\\
lo que se pierde   & interacciones entre nodos & nada de la rejilla\\
lo que se gana     & tres órdenes de magnitud menos de comparaciones & la interacción, si existe\\
\bottomrule
\end{tabularx}
\end{center}

\subsubsection*{El número de niveles que se pueden comparar a la vez}

Y aquí está el resultado que más conviene que el centro lea, porque va en contra de
la intuición de que más máquinas permiten comparar todo con todo.

Comparar todos los pares de 352 arms son $\binom{352}{2} = 61{,}776$ comparaciones.
Al coste medido de 0.6 minutos cada una son \textbf{798 horas}, que es más que el
coste de producir los 352 arms. Pero el problema no es el tiempo: es que 61{,}776
contrastes sobre nueve paneles son más de setecientas mil pruebas, y con esa
multiplicidad el máximo de una colección de diferencias ruidosas es grande aunque
todas las diferencias verdaderas sean cero.

\begin{aviso}{La cascada no es una concesión al presupuesto, es una necesidad estadística}
Con cómputo infinito se seguiría comparando contra un incumbente declarado y no
todos contra todos, porque el problema de la comparación múltiple no lo resuelve el
cómputo. Lo que cambia con la escala es \emph{cuántos arms se producen} y por tanto
cuántos candidatos entran en la cascada, no la forma de la comparación. El diseño
es el mismo a las dos escalas; a escala grande simplemente entra más gente al
torneo.
\end{aviso}

\subsection{Resumen del reparto}

\begin{center}
\begin{tabularx}{\linewidth}{@{}X c c@{}}
\toprule
 & acotada & a escala\\
\midrule
Nodos de decisión                & 10 & 10\\
Definiciones y marco             & idénticos & idénticos\\
Paneles                          & 9, sin agregar & 9, sin agregar\\
Suelo de población por celda     & 332 & 332\\
Estadístico y regla de cierre    & idénticos & idénticos\\
\midrule
Arms de predicción               & 29 & 352\\
Exploración                      & cascada & cruce\\
Interacciones entre nodos        & no observables & observables\\
Forma de la comparación          & contra incumbente & contra incumbente\\
Bancos de donantes igualados     & pendiente & sí\\
Escalera de recencia del banco   & 2 de 13 releases & 13 de 13\\
Encaminamiento con ajuste cruzado& pendiente & sí\\
\bottomrule
\end{tabularx}
\end{center}
